%------------------------------------------------------------------------------%
% Luis A. D'Afonseca
%
%------------------------------------------------------------------------------%
\documentclass[a4paper,12pt,fleqn]{article}

\usepackage{style_avaliacao}
\usepackage{systeme}

%\usepackage{comment}
%\excludecomment{resposta}
\newenvironment{resposta}%
  {\begingroup\color{blue}\small\begin{spacing}{0.9}}%
  {\end{spacing}\endgroup}

%------------------------------------------------------------------------------%
\begin{document}

\makeHeader{Integrais e Séries}{Prova 2 -- Turma 5}{03/10/2023}

% 01
%------------------------------------------------------------------------------%
\vspace{\baselineskip}
\exercise{20}
Use \textbf{integral por partes} para calcular a integral
\(
  \int_1^{e^2} \frac{\ln(x)}{\sqrt{x}} dx
\)
\begin{resposta}
  Calculando a integral indefinida por partes \(\int udv = uv - \int vdu\)

  Escolhendo
  \[
    u=\ln(x) \qquad dv = \frac{dx}{\sqrt{x}} = x^{-\sfrac{1}{2}} dx
  \]
  temos
  \[
    du = \frac{dx}{x}
    \qquad
    v = \frac{x^{\sfrac{1}{2}}}{\sfrac{1}{2}} = 2\sqrt{x}
  \]
  então
  \begin{align*}
    \int \frac{\ln(x)}{\sqrt{x}} \, dx
    & = \ln(x)2\sqrt{x} - \int 2\sqrt{x} \frac{dx}{x} \\
    & = 2\sqrt{x}\ln(x) - 2\int x^{\sfrac{1}{2}-1} dx \\
    & = 2\sqrt{x}\ln(x) - 2\int x^{-\sfrac{1}{2}} dx \\
    & = 2\sqrt{x}\ln(x) - 2\cdot 2\sqrt{x} + c \\[1mm]
    & = 2\sqrt{x} \left(\ln(x) - 2 \right) + c
  \end{align*}
  Calculando a integral definida
  \begin{align*}
  \int_1^{e^2} \frac{\ln(x)}{\sqrt{x}} dx
  & = 2\sqrt{x} \left(\ln(x) - 2 \right) \evalat{1}{e^2} \\
  & = \left[2\sqrt{e^2} \left(\ln\left(e^2\right) - 2 \right)\right]
    - \left[2\sqrt{1}   \left(\ln(1)              - 2 \right)\right] \\[1mm]
  & = \left[2e(2-2)\right] - \left[(2(-2))\right] \\
  & = 4
  \end{align*}

\end{resposta}

% 02
%------------------------------------------------------------------------------%
\exercise{20}
Use \textbf{substituição simples} para calcular a integral
\(
  \int \frac{a+bx^2}{\sqrt{3ax+bx^3}} dx
\)

\begin{resposta}
  Escolhendo a substituição
  \[
    u = 3ax+bx^2
    \qquad
    du = \left(3a + 3bx^2\right) dx = 3\left(a + bx^2\right) dx
  \]
  portanto
  \[
    \left(a+bx^2\right)dx = \frac{du}{3}
  \]
  então
  \begin{align*}
    \int \frac{a+bx^2}{\sqrt{3ax+bx^3}} dx
    & = \int \frac{1}{\sqrt{3ax+bx^3}} \left(a+bx^2\right)dx \\
    & = \int \frac{1}{\sqrt{u}} \frac{du}{3} \\
    & = \frac{1}{3} \int u^{-\sfrac{1}{2}} du \\
    & = \frac{1}{3} \frac{u^{\sfrac{1}{2}}}{\sfrac{1}{2}} + c \\
    & = \frac{2}{3} \sqrt{u} + c \\
    & = \frac{2}{3} \sqrt{3ax+bx^2} + c
  \end{align*}
\end{resposta}

% 03
%------------------------------------------------------------------------------%
\vspace{\baselineskip}
\exercise{40}
Escolha dois itens e calcule as primitivas das funções correspondentes
\begin{enumerate}[label=\alph*) \raisebox{-2mm}{\huge\(\square\)}]
  \setlength\itemsep{1em}
  \item \(\quad f(x) = \cos^2(x)\sin(2x)  \)
  \item \(\quad g(x) = \tan^4(x)\sec^6(x) \)
  \item \(\quad h(x) = \frac{x^5}{\sqrt{4-x^2}}  \)
\end{enumerate}
\textbf{Marque sua escolha nas caixas,
  caso contrário, serão corrigidos dois itens aleatoriamente.}

\begin{resposta}
  \begin{enumerate}[label=\alph*)]

    %--------------------------------------------------------------------------%
    \item
      \begin{align*}
        F
        & = \int f(x) dx \\
        & = \int \cos^2(x)\sin(2x) dx \\
        & = \int \cos^2(x)2\sin(x)\cos(x) dx \\
        & = 2\int \cos^3(x)\sin(x) dx
      \end{align*}
      Fazendo a substituição
      \(
        \qquad
        u = \cos(x)
        \qquad
        du = -\sin(x)dx
      \)
      \begin{align*}
        F
        & = -2\int u^3 du \\
        & = -2 \frac{u^4}{4} + c \\
        & = -\frac{1}{2}\cos^4(x) + c
      \end{align*}

    %--------------------------------------------------------------------------%
    \item
      \begin{align*}
        G
        & = \int \tan^4(x)\sec^6(x) dx \\
        & = \int \tan^4(x)\sec^4(x) \sec^2(x) dx \\
        & = \int \tan^4(x) \big(\sec^2(x)\big)^2 \sec^2(x) dx \\
        & = \int \tan^4(x) \big(\tan^2(x)+1\big)^2 \sec^2(x) dx
      \end{align*}
      Fazendo a substituição
      \(
        \qquad
        u = \tan(x)
        \qquad
        du = \sec^2(x) dx
      \)
      \begin{align*}
        G
        & = \int u^4 \big(u^2+1\big)^2 du \\
        & = \int u^4 \big(u^4+2u^2+1\big) du \\
        & = \int u^8 + 2u^6 + u^4 \,du \\
        & = \frac{u^9}{9} + 2\frac{u^7}{7} + \frac{u^5}{5} + c \\
        & = \frac{1}{9}\tan^9(x) + \frac{2}{7}\tan^7(x) + \frac{1}{5}\tan^5(x) + c
      \end{align*}
     %--------------------------------------------------------------------------%
     \item
      \[
        H = \int \frac{x^5}{\sqrt{4-x^2}} dx
      \]
      Fazendo a substituição
      \(
        x = 2\sin(\theta)
        \qquad
        dx = 2\cos(\theta) d\theta
      \)
      \begin{spreadlines}{2ex}
      \begin{align*}
        H
        & = \int \frac{2^5 \sin^5(\theta)}{\sqrt{4-4\sin^2(\theta)}} 2\cos(\theta) d\theta \\
        & = 2^5 \int \frac{\sin^5(\theta)}{\sqrt{1-\sin^2(\theta)}} \cos(\theta) d\theta \\
        & = 2^5 \int \frac{\sin^5(\theta)}{\cos(\theta)} \cos(\theta) d\theta \\
        & = 2^5 \int \sin^5(\theta) d\theta \\
        & = 2^5 \int \left(\sin^2(\theta)\right)^2 \sin(\theta) d\theta \\
        & = 2^5 \int \left(1-\cos^2(\theta)\right)^2 \sin(\theta) d\theta
      \end{align*}
    \end{spreadlines}
      Fazendo a substituição
      \(
        \qquad
        u = \cos(\theta)
        \qquad
        du = -\sin(\theta)d\theta
      \)
      \begin{align*}
        H
        & = -2^5 \int \left(1-u^2\right)^2 du \\
        & = -2^5 \int 1-2u^2 + u^4\, du \\
        & = 2^5 \int 2u^2 - u^4 -1\, du \\
        & = 2^5 \left(\frac{2}{3}u^3 - \frac{1}{5}u^5 - u \right) + c \\
        & = 2^5 \left(\frac{2}{3}\cos^3(\theta)
         - \frac{1}{5}\cos^5(\theta)
         - \cos(\theta) \right) + c
      \end{align*}
      Para calcular \(\cos(\theta)\) usamos que \(\sin(\theta) = \sfrac{x}{2}\), portanto
      a hipotenusa é 2 e o cateto oposto é $x$ assim o cateto adjacente é
      \(
        a = \sqrt{2^2-x^2} = \sqrt{4-x^2}
      \)
      e o cosseno é
      \[
        \cos(\theta) = \frac{\sqrt{4-x^2}}{2}
      \]
      Voltando para a integral
      \begin{spreadlines}{2ex}
      \begin{align*}
        H
        & = 2^5\left[
            \frac{2}{3}\left(\frac{\sqrt{4-x^2}}{2}\right)^3
            - \frac{1}{5}\left(\frac{\sqrt{4-x^2}}{2}\right)^5
            - \frac{\sqrt{4-x^2}}{2}
          \right] + c \\
        & = 2^5\left[
            \frac{2}{3\cdot 2^3}\left(\sqrt{4-x^2}\right)^3
          - \frac{1}{5\cdot 2^5}\left(\sqrt{4-x^2}\right)^5
          - \frac{1}{2} \sqrt{4-x^2}
          \right] + c \\
        & = \sqrt{4-x^2}
            \left[
              \frac{2^3}{3}\left(\sqrt{4-x^2}\right)^2
            - \frac{1}  {5}\left(\sqrt{4-x^2}\right)^4
            - 2^4
            \right] + c \\
        & = \sqrt{4-x^2}
            \left[
              \frac{8}{3}\left(4-x^2\right)
            - \frac{1}{5}\left(4-x^2\right)^2
            - 16
            \right] + c \\
        & = \sqrt{4-x^2}
            \left[
              \frac{32-8x^2}{3}
            - \frac{16-8x^2+x^4}{5}
            - 16
            \right] + c \\
        & = \sqrt{4-x^2}
            \left[
               \frac{5\cdot 32 - 5\cdot 8x^2}{15}
              -\frac{3\cdot 16 - 3\cdot 8x^2 + 3x^4}{15} - \frac{16\cdot 15}{15}
            \right] + c \\
        & = \frac{\sqrt{4-x^2}}{15}
            \Big[
              16\cdot 5\cdot 2 - 16 \cdot 3 - 16\cdot 5 \cdot 3
            - 5\cdot 8x^2 + 3\cdot 8x^2
            - 3x^4
            \Big] + c \\
        & = \frac{\sqrt{4-x^2}}{15}
            \Big[
            - 16\cdot 5 -16 \cdot 3
            - 2\cdot 8x^2
            - 3x^4
            \Big] + c \\
        & = -\frac{\sqrt{4-x^2}}{15}
            \Big[
              16\cdot 8
            + 2\cdot 8x^2
            + 3x^4
            \Big] + c \\
        & = -\frac{\sqrt{4-x^2}}{15}
            \Big[ 128 + 16x^2 + 3x^4 \Big] + c
      \end{align*}
      \end{spreadlines}

%      \[
%      H = \int \frac{dx}{x^5\sqrt{9x^2-1}}
%      \]
%      Fazemos a substituição
%      \(
%      x = \frac{1}{3}\sec(\theta)
%      \qquad
%      dx = \frac{1}{3}\tan(\theta)\sec{\theta} d\theta
%      \)
%      \begin{align*}
%        H
%        & = \int \frac{\sfrac{1}{3}\tan(\theta)\sec(\theta)}
%        {\left(\sfrac{1}{3}\sec(\theta)\right)^5
%          \sqrt{9\left(\sfrac{1}{3}\sec(\theta)\right)^2-1}} d\theta \\
%        & = \int \frac{\sfrac{1}{3}\tan(\theta)\sec(\theta)}
%        {\sfrac{1}{3^5}\sec^5(\theta)
%          \sqrt{\sec^2(\theta)-1}} d\theta \\
%        & = 3^4 \int \frac{\tan(\theta)\sec(\theta)}
%        {\sec^5(\theta)\tan(\theta)} d\theta \\
%        & = 3^4 \int \frac{1}{\sec^4(\theta)} d\theta \\
%        & = 3^4 \int \frac{\sec^2(\theta)}{\sec^6(\theta)} d\theta \\
%        & = 3^4 \int \frac{\sec^2(\theta)}{\left(\sec^2(\theta)\right)^3}d\theta \\
%        & = 3^4 \int \frac{\sec^2(\theta)}{\left(\tan^2(\theta)+1\right)^3}d\theta
%      \end{align*}
%      Fazemos a substituição \(u=\sec(\theta)\)
%
  \end{enumerate}
\end{resposta}

% 04
%------------------------------------------------------------------------------%
\vspace{\baselineskip}
\exercise{20}
Calcule o volume do sólido de revolução gerado pela rotação,
em torno da reta \(x=-1\), da região contida entre as curvas
\[
  f(x) = \frac{1}{4x-x^2} \qquad
  y = 0 \qquad
  x = 1 \qquad
  x = 3
\]

\begin{resposta}
  Volume por castas cilíndricas
  \begin{align*}
    V & = \int_a^b 2\pi r(x) h(x) dx \\
    a & = 1 \\
    b & = 3 \\
    r & = x+1 \\
    h & = \frac{1}{4x-x^2}
  \end{align*}
  portanto
  \[
    V = 2\pi \int_1^3 \frac{x+1}{4x-x^2} dx
  \]
  Calculando a integral indefinida
  \[
    F = \int \frac{x+1}{4x-x^2} dx = \int \frac{x+1}{x(4-x)} dx
  \]
  Por frações parciais temos
  \begin{align*}
    \frac{x+1}{x(4-x)} & = \frac{A}{x} + \frac{B}{4-x} \\
    x+1                & = A(4-x) + Bx                 \\
                       & = (B-A)x + 4A
  \end{align*}
  Igualando os coeficientes
  \[
    \systeme[AB]{B-A=1,4A=1}
  \]
  obtemos os valores
  \(A=\frac{1}{4}\) e
  \(B=\frac{5}{4}\)
  portanto
  \begin{align*}
    F & = \int \frac{1}{4x} + \frac{5}{4(4-x)} dx \\
      & = \frac{1}{4}\ln(x) - \frac{5}{4}\ln(4-x) + c
  \end{align*}
  Voltando ao volume temos
  \begin{align*}
    V & = 2\pi F(x)\evalat{1}{3} \\
      & = 2\pi \left[\frac{1}{4}\ln(x) - \frac{5}{4}\ln(4-x)\right] \evalat{1}{3} \\
      & = 2\pi \left[
                 \left(\frac{1}{4}\ln(3) - \frac{5}{4}\ln(4-3)\right)
               - \left(\frac{1}{4}\ln(1) - \frac{5}{4}\ln(4-1)\right)
               \right] \\
      & = 2\pi \left(\frac{1}{4}\ln(3) + \frac{5}{4}\ln(3)\right) \\
      & = 2\pi \frac{6}{4}\ln(3) \\
      & = 3\pi \ln(3)
  \end{align*}
\end{resposta}

%------------------------------------------------------------------------------%
\end{document}
%------------------------------------------------------------------------------%
